\section{Strategi Optimasi pada Level Instruksi}
\label{sec:optimasi-instruksi}

Setelah hambatan ditemukan melalui pengukuran, langkah berikutnya
adalah memperbaiki instruksinya. Bagian ini membahas enam teknik
optimasi yang dapat diterapkan langsung pada program assembly 8086,
kemudian menutup dengan trade-off, cara pengukuran yang adil, dan
batas etika optimasi.

\subsection{Memilih Register daripada Memori}

Teknik paling mendasar adalah menyimpan nilai yang sering dipakai di
dalam register, bukan di memori. Perhatikan perbedaannya:

\begin{lstlisting}[language=MASM, numbers=none]
; --- kurang baik: nilai diambil dari memori setiap iterasi ---
    MOV CX, 10
ULANG:
    ADD AX, jumlah      ; akses memori: 9 + 6 = 15 siklus
    LOOP ULANG

; --- lebih baik: nilai sudah ada di register ---
    MOV CX, 10
    MOV BX, jumlah      ; sekali saja: 8 + 6 = 14 siklus
ULANG:
    ADD AX, BX          ; akses register: 3 siklus
    LOOP ULANG
\end{lstlisting}

Pada versi pertama, sepuluh iterasi memerlukan sepuluh akses memori
($10 \times 15 = 150$ siklus). Pada versi kedua, satu akses memori di
awal ditambah sepuluh operasi register
($14 + 10 \times 3 = 44$ siklus). Selisih ini berasal dari satu
perubahan sederhana: memindahkan data ke tempat yang paling cepat
diakses prosesor.

\subsection{Strength Reduction}

\term{Strength reduction} adalah penggantian operasi yang mahal
dengan operasi yang lebih murah yang menghasilkan nilai sama.
Tabel~\ref{tab:strength-reduction} memuat penggantian yang paling
umum beserta syarat berlakunya. Kolom syarat sangat penting: tanpa
memenuhi syarat tersebut, penggantian justru menghasilkan nilai yang
salah.

\begin{table}[H]
    \centering
    \caption{Penggantian operasi (\eng{strength reduction}) yang umum
    pada 8086.}
    \label{tab:strength-reduction}
    \small
    \begin{tabular}{@{}p{3.8cm}p{3.8cm}p{5.3cm}@{}}
        \toprule
        \textbf{Operasi asal} & \textbf{Pengganti} &
        \textbf{Syarat berlaku} \\
        \midrule
        \lstasm|MUL BX| dengan \reg{BX} = 2 & \lstasm|ADD AX, AX| &
        Pengali tepat dua; hasil tidak memerlukan \reg{DX} \\
        \addlinespace
        \lstasm|MUL BX| dengan \reg{BX} = $2^n$ & \lstasm|SHL AX, n| &
        Pengali pangkat dua dan tak bertanda \\
        \addlinespace
        \lstasm|DIV BX| dengan \reg{BX} = $2^n$ & \lstasm|SHR AX, n| &
        Pembagi pangkat dua, nilai tak bertanda,
        sisa pembagian tidak diperlukan \\
        \addlinespace
        \lstasm|IDIV BX| dengan \reg{BX} = $2^n$ & \lstasm|SAR AX, n| &
        Nilai bertanda (komplemen dua); pembulatan ke bawah \\
        \addlinespace
        \lstasm|MOV AX, 0| & \lstasm|XOR AX, AX| &
        Hasil harus nol; \eng{flag} memang boleh berubah \\
        \addlinespace
        \lstasm|CMP AX, 0| & \lstasm|TEST AX, AX| &
        Uji dilakukan hanya terhadap nol \\
        \addlinespace
        \lstasm|ADD AX, 1| & \lstasm|INC AX| &
        \flag{CF} tidak diperiksa sesudahnya (\asm{INC} tidak mengubah
        \flag{CF}) \\
        \addlinespace
        \lstasm|SUB AX, 1| & \lstasm|DEC AX| &
        \flag{CF} tidak diperiksa sesudahnya \\
        \addlinespace
        \lstasm|ADD SI, 2| & \lstasm|LEA SI, [SI+2]| &
        \eng{Flag} harus dipertahankan agar pembandingan sebelumnya
        tetap sahih \\
        \addlinespace
        \lstasm|DEC CX| + \lstasm|JNZ| & \lstasm|LOOP| &
        Penghitung ada di \reg{CX} dan jarak lompatan kurang dari
        $\pm$128 byte \\
        \bottomrule
    \end{tabular}
\end{table}

\subsection{Pergeseran sebagai Pengganti Perkalian dan Pembagian}

Penggantian \asm{MUL} dan \asm{DIV} dengan pergeseran adalah contoh
\eng{strength reduction} yang paling menguntungkan pada 8086.
Sebuah \lstasm|MUL BX| memerlukan 118--133 siklus, sedangkan
\lstasm|SHL AX, 1| hanya 2 siklus. Untuk mengalikan dengan delapan,
tiga kali pergeseran satu bit sudah cukup dan totalnya hanya enam
siklus.

Menariknya, pada 8086 ada dua cara yang keduanya sah, tetapi
memiliki keunggulan berbeda:

\begin{lstlisting}[language=MASM, numbers=none]
; --- varian cepat: 3 instruksi, 6 byte, 6 siklus ---
    SHL AX, 1           ; AX = AX x 2
    SHL AX, 1           ; AX = AX x 4
    SHL AX, 1           ; AX = AX x 8

; --- varian kecil: 2 instruksi, 4 byte, 20 siklus ---
    MOV CL, 3           ; CL = banyak pergeseran
    SHL AX, CL          ; 8 + 4 x 3 = 20 siklus
\end{lstlisting}

Varian pertama lebih cepat tetapi lebih panjang dua byte; varian
kedua lebih pendek tetapi lebih lambat. Contoh ini menunjukkan bahwa
``teroptimasi'' tidak memiliki satu arti tunggal: jawabannya
bergantung pada metrik yang dipentingkan.

\begin{peringatanbox}
Pergeseran hanya setara dengan perkalian bila nilai yang digeser
\emph{tidak bertanda} dan bit yang tergeser keluar memang boleh
hilang. Untuk nilai bertanda, \asm{SHR} dapat mengubah tanda menjadi
positif; pada kasus itu \asm{SAR} yang benar. Selain itu, \asm{MUL}
memperhatikan \flag{CF} dan \flag{OF} untuk mendeteksi limpahan
(\eng{overflow}), sedangkan \asm{SHL} tidak. Bila program mengandalkan
kedua \eng{flag} tersebut, penggantian harus dibatalkan.
\end{peringatanbox}

\subsection{LEA untuk Aritmatika Tanpa Mengubah Flag}

Instruksi \asm{LEA} (\eng{Load Effective Address}) menghitung alamat
efektif tanpa benar-benar mengakses memori. Instruksi ini memiliki
dua kegunaan yang sering terlewat.

\begin{enumerate}
    \item \textbf{Menghitung alamat majemuk dalam satu instruksi.}
          \lstasm|LEA SI, [BX+SI+4]| menjumlahkan \reg{BX}, \reg{SI},
          dan konstanta 4 sekaligus, sesuatu yang tidak dapat
          dilakukan satu instruksi aritmatika biasa.
    \item \textbf{Aritmatika tanpa mengubah flag.} Karena \asm{LEA}
          tidak menyentuh \eng{flag}, ia dipakai ketika sebuah
          pembandingan harus tetap berlaku setelah alamat dihitung.
\end{enumerate}

Perhatikan ilustrasi berikut. Pada versi pertama, \asm{ADD} mengubah
\flag{ZF} sehingga lompatan bersyarat menjadi salah:

\begin{lstlisting}[language=MASM, numbers=none]
    CMP AX, 0           ; ZF = 1 bila AX = 0
    ADD SI, 2           ; SALAH: ADD mengubah ZF
    JE  kosong          ; lompatan tidak lagi dapat dipercaya

    CMP AX, 0           ; ZF = 1 bila AX = 0
    LEA SI, [SI+2]      ; BENAR: LEA tidak mengubah flag
    JE  kosong          ; lompatan masih sahih
\end{lstlisting}

Namun \asm{LEA} bukan pilihan yang selalu lebih cepat: pada 8086
\lstasm|LEA SI,[SI+2]| memerlukan 10 siklus, sedangkan
\lstasm|ADD SI,2| hanya 4 siklus. Keunggulan \asm{LEA} adalah
keringkasan dan sifatnya yang tidak mengubah \eng{flag}, bukan
kecepatannya. Hubungan \asm{LEA} dengan mode pengalamatan dibahas
pada Bab~\ref{chap:mode-pengalamatan}.

\subsection{Mengurangi Lompatan dan \eng{Loop Unrolling}}

Setiap lompatan memiliki harga: \lstasm|LOOP| memerlukan 17 siklus,
lompatan bersyarat 16 siklus, dan lompatan jauh lebih mahal lagi.
Bahkan pada prosesor 8086, lompatan juga menghalangi
\eng{prefetch queue} mengisi instruksi berikutnya. Ada dua cara
mengurangi jumlah lompatan.

\begin{enumerate}
    \item \textbf{Menyusun ulang alur.} Bila dua cabang berakhir pada
          pekerjaan yang sama, pekerjaan itu dapat dipindahkan keluar
          dari cabang sehingga hanya ada satu salinan. Susunan
          ``lompat ke pekerjaan bersama'' menggantikan dua salinan
          kode menjadi satu.
    \item \textbf{\eng{Loop unrolling}.} Isi perulangan digandakan
          sehingga satu iterasi mengerjakan dua elemen sekaligus.
          Penghitung dan lompatan dikerjakan setengah kali, sehingga
          overhead \asm{LOOP} berkurang separuh. Harganya jelas:
          kode menjadi lebih panjang.
\end{enumerate}

\begin{lstlisting}[language=MASM, numbers=none]
; --- perulangan biasa: 1 elemen per iterasi ---
    MOV CX, 10
ULANG:
    ADD AX, [SI]
    ADD SI, 2
    LOOP ULANG

; --- setelah unrolling 2x: 2 elemen per iterasi ---
    MOV CX, 5           ; cukup 5 iterasi
ULANG2:
    ADD AX, [SI]        ; elemen ke-1
    ADD AX, [SI+2]      ; elemen ke-2
    ADD SI, 4           ; maju dua word
    LOOP ULANG2
\end{lstlisting}

\subsection{Mengosongkan Register dengan \asm{XOR reg, reg}}
\label{sec:xor-reg-reg}

Untuk mengosongkan sebuah register, banyak programmer menulis
\lstasm|MOV AX, 0|. Tulisan itu benar dan mudah dibaca, tetapi bukan
yang terpendek:

\begin{lstlisting}[language=MASM, numbers=none]
    MOV AX, 0           ; 3 byte (B8 00 00), 4 siklus
    XOR AX, AX          ; 2 byte (31 C0),    3 siklus
    SUB AX, AX          ; 2 byte (29 C0),    3 siklus
\end{lstlisting}

\lstasm|XOR AX, AX| lebih pendek satu byte dan lebih cepat satu
siklus, karena tidak perlu memuat nilai segera. Selain itu,
\lstasm|XOR AX, AX| mengosongkan register 16-bit sekaligus, sedangkan
\lstasm|MOV AX, 0| menyimpan nilai \hex{00} dua byte. Dasar logika
operasi \eng{XOR} dijelaskan pada Bab~\ref{chap:logika}.

Harganya adalah \eng{flag}: \asm{XOR} mengubah \flag{ZF}, \flag{SF},
\flag{PF}, dan menghapus \flag{CF} serta \flag{OF}. Bila
pengosongan register dilakukan di antara sebuah pembandingan dan
lompatan bersyaratnya, \asm{XOR} justru merusak program. Karena itu
pemakaiannya harus ditempatkan pada awal blok, bukan di tengah alur
keputusan.

\subsection{Trade-off: Ukuran Kode atau Kecepatan}

Tidak ada optimasi yang gratis. Setiap perubahan hampir selalu
memindahkan biaya dari satu metrik ke metrik lain. Dua arah optimasi
yang berlawanan dirumuskan berikut ini.

\begin{itemize}
    \item \textbf{Optimasi ukuran} dipilih bila kode harus muat pada
          memori terbatas, misalnya \eng{firmware}, rutin
          \eng{bootloader}, atau program yang dipanggil dari disket.
          Tekniknya: varian pendek (\lstasm|MOV CL,3| +
          \lstasm|SHL AX, CL|), hindari \eng{unrolling}, dan hindari
          salinan kode yang berulang.
    \item \textbf{Optimasi kecepatan} dipilih bila program berjalan
          pada bagian yang paling sering dieksekusi dan memori masih
          mencukupi. Tekniknya: varian cepat (tiga kali
          \lstasm|SHL AX,1|), \eng{unrolling}, register daripada
          memori, dan \asm{LEA}.
\end{itemize}

Gambar~\ref{fig:alur-keputusan-optimasi} merangkum alur keputusan yang
dipakai pada praktikum akhir bab ini.

\begin{figure}[H]
    \centering
    \begin{tikzpicture}[
        node distance=0.6cm,
        simpul/.style={
            draw=obeBlue, line width=0.8pt, fill=obeBlueLight,
            rounded corners=3pt, align=center, font=\small,
            inner sep=5pt, text width=4.3cm
        },
        uji/.style={
            draw=obeAmber, line width=0.8pt, fill=obeAmberLight,
            diamond, aspect=2.6, align=center, font=\small,
            inner sep=1pt, text width=3.2cm
        },
        akhir/.style={
            draw=obeGreen, line width=0.8pt, fill=obeGreenLight,
            rounded corners=3pt, align=center, font=\small,
            inner sep=5pt, text width=4.3cm
        },
        panah/.style={-{Stealth[length=2.2mm]}, line width=0.8pt, obeTeal},
        ket/.style={font=\scriptsize, text=obeGray, fill=white,
                    inner sep=1pt}
    ]
        \node[simpul] (ukur1) {Langkah 1: ukur versi naif (waktu,
            ukuran byte)};
        \node[uji, below=of ukur1] (tanya)
            {\textbf{Hambatan utama sudah diketahui?}};
        \node[simpul, below=of tanya, fill=obeMaroonLight,
              draw=obeMaroon] (cari)
            {Langkah 2: cari hambatan dengan \code{U} dan penghitung
             iterasi};
        \node[uji, below=of cari] (pilih)
            {\textbf{Batasan pokok proyek?}};
        \node[akhir, below left=0.8cm and 0.15cm of pilih] (kecil)
            {Optimasi ukuran: varian pendek, tanpa \eng{unrolling}};
        \node[akhir, below right=0.8cm and 0.15cm of pilih] (cepat)
            {Optimasi kecepatan: register, tiga kali \asm{SHL AX,1},
             \eng{unrolling}};
        \node[akhir, below=2.5cm of pilih] (verif)
            {Langkah 3: ukur ulang dan verifikasi hasil tetap benar};

        \draw[panah] (ukur1) -- (tanya);
        \draw[panah] (tanya) -- node[ket, right] {belum} (cari);
        \draw[panah] (cari.west) -- ++(-0.9,0) |-
            node[ket, left, pos=0.3] {ulangi} (tanya.west);
        \draw[panah] (tanya) -- (pilih);
        \draw[panah] (pilih) -- node[ket, left] {ukuran} (kecil);
        \draw[panah] (pilih) -- node[ket, right] {kecepatan} (cepat);
        \draw[panah] (kecil.south) |- (verif.west);
        \draw[panah] (cepat.south) |- (verif.east);
    \end{tikzpicture}
    \caption{Alur keputusan optimasi: ukur dahulu, kenali hambatan,
    pilih arah optimasi, lalu ukur ulang.}
    \label{fig:alur-keputusan-optimasi}
\end{figure}

\subsection{Pengukuran yang Adil}

Agar perbandingan sebelum dan sesudah dapat dipercaya, tiga syarat
berikut harus dipenuhi.

\begin{enumerate}
    \item \textbf{Masukan sama dan jumlah iterasi sama.} Kedua versi
          harus memakai larik yang sama dan mengerjakan pekerjaan yang
          sama banyaknya. Penghitung iterasi dicatat untuk
          membuktikannya.
    \item \textbf{Batas pengukuran sama.} Waktu mulai dicatat tepat
          sebelum pekerjaan dan waktu berhenti tepat sesudah
          pekerjaan, tidak termasuk pemuatan program, pencetakan
          keluaran, atau penekanan tombol.
    \item \textbf{Diulang beberapa kali.} Jalankan minimal lima kali,
          buang percobaan pertama (pemanasan), lalu ambil nilai
          tengah. Satu kali pengukuran selalu ramai oleh kebetulan.
\end{enumerate}

Waktu dapat dicatat dengan layanan sistem. BIOS menyediakan pencacah
\eng{tick} melalui \code{INT 1Ah} fungsi \hex{00}, yang bertambah
18,2 kali per detik.

\begin{lstlisting}[language=MASM, numbers=none]
    MOV AH, 00h         ; baca pencacah tick BIOS
    INT 1Ah             ; CX:DX = tick sejak tengah malam
    MOV tick_awal, DX   ; simpan titik awal

    ; ... pekerjaan yang diukur ...

    MOV AH, 00h
    INT 1Ah             ; baca lagi
    SUB DX, tick_awal   ; DX = selisih tick (18,2 tick = 1 detik)
\end{lstlisting}

Cara lain adalah \code{INT 21h} fungsi \hex{2C} yang mengembalikan
jam, menit, detik, dan perseratus detik pada register \reg{CH},
\reg{CL}, \reg{DH}, dan \reg{DL}. Karena resolusi pencacah
\eng{tick} hanya sekitar 55 milidetik, pekerjaan yang diukur harus
dijalankan berulang-ulang di dalam satu perulangan pengukuran agar
selisihnya terlihat.

\subsection{Etika dan Batas Optimasi}

Optimasi memiliki harga yang tidak terlihat pada metrik: keterbacaan.
Kode yang penuh trik sulit dibaca, sulit diperbaiki, dan mudah rusak
ketika kebutuhan berubah. Tiga prinsip berikut menjadi pedoman.

\begin{itemize}
    \item \textbf{Benar lebih dahulu, cepat kemudian.} Program yang
          cepat tetapi salah tidak berguna. Uji kebenaran setiap versi
          dengan masukan yang sama sebelum membandingkan kecepatannya.
    \item \textbf{Jangan mengoptimasi dini.} Optimasi hanya layak
          dilakukan pada bagian yang sudah terbukti menjadi hambatan
          melalui pengukuran, bukan pada bagian yang hanya
          \emph{diperkirakan} lambat.
    \item \textbf{Jaga keterbacaan dan dokumentasi.} Setiap trik harus
          diberi komentar yang menjelaskan alasan dan syarat
          berlakunya. Simpan versi naif sebagai pembanding dan
          rujukan bila versi teroptimasi bermasalah.
\end{itemize}

\begin{infobox}
Optimasi yang baik bersifat dapat dipertanggungjawabkan: ada
pengukuran sebelum, ada penjelasan mengapa perubahan dipilih, dan ada
pengukuran sesudah. Bila salah satu dari ketiganya tidak ada, yang
dikerjakan bukan optimasi melainkan sekadar perubahan.
\end{infobox}
